% Options for packages loaded elsewhere
\PassOptionsToPackage{unicode}{hyperref}
\PassOptionsToPackage{hyphens}{url}
%
\documentclass[
]{article}
\usepackage{amsmath,amssymb}
\usepackage{iftex}
\ifPDFTeX
  \usepackage[T1]{fontenc}
  \usepackage[utf8]{inputenc}
  \usepackage{textcomp} % provide euro and other symbols
\else % if luatex or xetex
  \usepackage{unicode-math} % this also loads fontspec
  \defaultfontfeatures{Scale=MatchLowercase}
  \defaultfontfeatures[\rmfamily]{Ligatures=TeX,Scale=1}
\fi
\usepackage{lmodern}
\ifPDFTeX\else
  % xetex/luatex font selection
\fi
% Use upquote if available, for straight quotes in verbatim environments
\IfFileExists{upquote.sty}{\usepackage{upquote}}{}
\IfFileExists{microtype.sty}{% use microtype if available
  \usepackage[]{microtype}
  \UseMicrotypeSet[protrusion]{basicmath} % disable protrusion for tt fonts
}{}
\makeatletter
\@ifundefined{KOMAClassName}{% if non-KOMA class
  \IfFileExists{parskip.sty}{%
    \usepackage{parskip}
  }{% else
    \setlength{\parindent}{0pt}
    \setlength{\parskip}{6pt plus 2pt minus 1pt}}
}{% if KOMA class
  \KOMAoptions{parskip=half}}
\makeatother
\usepackage{xcolor}
\usepackage[margin=1in]{geometry}
\usepackage{color}
\usepackage{fancyvrb}
\newcommand{\VerbBar}{|}
\newcommand{\VERB}{\Verb[commandchars=\\\{\}]}
\DefineVerbatimEnvironment{Highlighting}{Verbatim}{commandchars=\\\{\}}
% Add ',fontsize=\small' for more characters per line
\usepackage{framed}
\definecolor{shadecolor}{RGB}{248,248,248}
\newenvironment{Shaded}{\begin{snugshade}}{\end{snugshade}}
\newcommand{\AlertTok}[1]{\textcolor[rgb]{0.94,0.16,0.16}{#1}}
\newcommand{\AnnotationTok}[1]{\textcolor[rgb]{0.56,0.35,0.01}{\textbf{\textit{#1}}}}
\newcommand{\AttributeTok}[1]{\textcolor[rgb]{0.13,0.29,0.53}{#1}}
\newcommand{\BaseNTok}[1]{\textcolor[rgb]{0.00,0.00,0.81}{#1}}
\newcommand{\BuiltInTok}[1]{#1}
\newcommand{\CharTok}[1]{\textcolor[rgb]{0.31,0.60,0.02}{#1}}
\newcommand{\CommentTok}[1]{\textcolor[rgb]{0.56,0.35,0.01}{\textit{#1}}}
\newcommand{\CommentVarTok}[1]{\textcolor[rgb]{0.56,0.35,0.01}{\textbf{\textit{#1}}}}
\newcommand{\ConstantTok}[1]{\textcolor[rgb]{0.56,0.35,0.01}{#1}}
\newcommand{\ControlFlowTok}[1]{\textcolor[rgb]{0.13,0.29,0.53}{\textbf{#1}}}
\newcommand{\DataTypeTok}[1]{\textcolor[rgb]{0.13,0.29,0.53}{#1}}
\newcommand{\DecValTok}[1]{\textcolor[rgb]{0.00,0.00,0.81}{#1}}
\newcommand{\DocumentationTok}[1]{\textcolor[rgb]{0.56,0.35,0.01}{\textbf{\textit{#1}}}}
\newcommand{\ErrorTok}[1]{\textcolor[rgb]{0.64,0.00,0.00}{\textbf{#1}}}
\newcommand{\ExtensionTok}[1]{#1}
\newcommand{\FloatTok}[1]{\textcolor[rgb]{0.00,0.00,0.81}{#1}}
\newcommand{\FunctionTok}[1]{\textcolor[rgb]{0.13,0.29,0.53}{\textbf{#1}}}
\newcommand{\ImportTok}[1]{#1}
\newcommand{\InformationTok}[1]{\textcolor[rgb]{0.56,0.35,0.01}{\textbf{\textit{#1}}}}
\newcommand{\KeywordTok}[1]{\textcolor[rgb]{0.13,0.29,0.53}{\textbf{#1}}}
\newcommand{\NormalTok}[1]{#1}
\newcommand{\OperatorTok}[1]{\textcolor[rgb]{0.81,0.36,0.00}{\textbf{#1}}}
\newcommand{\OtherTok}[1]{\textcolor[rgb]{0.56,0.35,0.01}{#1}}
\newcommand{\PreprocessorTok}[1]{\textcolor[rgb]{0.56,0.35,0.01}{\textit{#1}}}
\newcommand{\RegionMarkerTok}[1]{#1}
\newcommand{\SpecialCharTok}[1]{\textcolor[rgb]{0.81,0.36,0.00}{\textbf{#1}}}
\newcommand{\SpecialStringTok}[1]{\textcolor[rgb]{0.31,0.60,0.02}{#1}}
\newcommand{\StringTok}[1]{\textcolor[rgb]{0.31,0.60,0.02}{#1}}
\newcommand{\VariableTok}[1]{\textcolor[rgb]{0.00,0.00,0.00}{#1}}
\newcommand{\VerbatimStringTok}[1]{\textcolor[rgb]{0.31,0.60,0.02}{#1}}
\newcommand{\WarningTok}[1]{\textcolor[rgb]{0.56,0.35,0.01}{\textbf{\textit{#1}}}}
\usepackage{graphicx}
\makeatletter
\newsavebox\pandoc@box
\newcommand*\pandocbounded[1]{% scales image to fit in text height/width
  \sbox\pandoc@box{#1}%
  \Gscale@div\@tempa{\textheight}{\dimexpr\ht\pandoc@box+\dp\pandoc@box\relax}%
  \Gscale@div\@tempb{\linewidth}{\wd\pandoc@box}%
  \ifdim\@tempb\p@<\@tempa\p@\let\@tempa\@tempb\fi% select the smaller of both
  \ifdim\@tempa\p@<\p@\scalebox{\@tempa}{\usebox\pandoc@box}%
  \else\usebox{\pandoc@box}%
  \fi%
}
% Set default figure placement to htbp
\def\fps@figure{htbp}
\makeatother
\setlength{\emergencystretch}{3em} % prevent overfull lines
\providecommand{\tightlist}{%
  \setlength{\itemsep}{0pt}\setlength{\parskip}{0pt}}
\setcounter{secnumdepth}{-\maxdimen} % remove section numbering
\usepackage{bookmark}
\IfFileExists{xurl.sty}{\usepackage{xurl}}{} % add URL line breaks if available
\urlstyle{same}
\hypersetup{
  pdftitle={The normal distribution},
  hidelinks,
  pdfcreator={LaTeX via pandoc}}

\title{The normal distribution}
\author{}
\date{\vspace{-2.5em}}

\begin{document}
\maketitle

In this lab, you'll investigate the probability distribution that is
most central to statistics: the normal distribution. If you are
confident that your data are nearly normal, that opens the door to many
powerful statistical methods. Here we'll use the graphical tools of R to
assess the normality of our data and also learn how to generate random
numbers from a normal distribution.

\subsection{Getting Started}\label{getting-started}

\subsubsection{The data}\label{the-data}

This week you'll be working with fast food data. This data set contains
data on 515 menu items from some of the most popular fast food
restaurants worldwide. Let's take a quick peek at the first few rows of
the data.

Either you can use \texttt{glimpse} like before, or \texttt{head} to do
this.

\begin{Shaded}
\begin{Highlighting}[]
\FunctionTok{library}\NormalTok{(tidyverse)}
\FunctionTok{library}\NormalTok{(openintro)}
\FunctionTok{glimpse}\NormalTok{(fastfood)}
\end{Highlighting}
\end{Shaded}

\begin{verbatim}
## Rows: 515
## Columns: 17
## $ restaurant  <chr> "Mcdonalds", "Mcdonalds", "Mcdonalds", "Mcdonalds", "Mcdon~
## $ item        <chr> "Artisan Grilled Chicken Sandwich", "Single Bacon Smokehou~
## $ calories    <dbl> 380, 840, 1130, 750, 920, 540, 300, 510, 430, 770, 380, 62~
## $ cal_fat     <dbl> 60, 410, 600, 280, 410, 250, 100, 210, 190, 400, 170, 300,~
## $ total_fat   <dbl> 7, 45, 67, 31, 45, 28, 12, 24, 21, 45, 18, 34, 20, 34, 8, ~
## $ sat_fat     <dbl> 2.0, 17.0, 27.0, 10.0, 12.0, 10.0, 5.0, 4.0, 11.0, 21.0, 4~
## $ trans_fat   <dbl> 0.0, 1.5, 3.0, 0.5, 0.5, 1.0, 0.5, 0.0, 1.0, 2.5, 0.0, 1.5~
## $ cholesterol <dbl> 95, 130, 220, 155, 120, 80, 40, 65, 85, 175, 40, 95, 125, ~
## $ sodium      <dbl> 1110, 1580, 1920, 1940, 1980, 950, 680, 1040, 1040, 1290, ~
## $ total_carb  <dbl> 44, 62, 63, 62, 81, 46, 33, 49, 35, 42, 38, 48, 48, 67, 31~
## $ fiber       <dbl> 3, 2, 3, 2, 4, 3, 2, 3, 2, 3, 2, 3, 3, 5, 2, 2, 3, 3, 5, 2~
## $ sugar       <dbl> 11, 18, 18, 18, 18, 9, 7, 6, 7, 10, 5, 11, 11, 11, 6, 3, 1~
## $ protein     <dbl> 37, 46, 70, 55, 46, 25, 15, 25, 25, 51, 15, 32, 42, 33, 13~
## $ vit_a       <dbl> 4, 6, 10, 6, 6, 10, 10, 0, 20, 20, 2, 10, 10, 10, 2, 4, 6,~
## $ vit_c       <dbl> 20, 20, 20, 25, 20, 2, 2, 4, 4, 6, 0, 10, 20, 15, 2, 6, 15~
## $ calcium     <dbl> 20, 20, 50, 20, 20, 15, 10, 2, 15, 20, 15, 35, 35, 35, 4, ~
## $ salad       <chr> "Other", "Other", "Other", "Other", "Other", "Other", "Oth~
\end{verbatim}

\begin{Shaded}
\begin{Highlighting}[]
\FunctionTok{table}\NormalTok{(fastfood}\SpecialCharTok{$}\NormalTok{restaurant)}
\end{Highlighting}
\end{Shaded}

\begin{verbatim}
## 
##       Arbys Burger King Chick Fil-A Dairy Queen   Mcdonalds       Sonic 
##          55          70          27          42          57          53 
##      Subway   Taco Bell 
##          96         115
\end{verbatim}

You'll see that for every observation there are 17 measurements, many of
which are nutritional facts.

You'll be focusing on just three columns to get started: restaurant,
calories, calories from fat.

Let's first focus on just products from McDonalds and Dairy Queen.

\begin{Shaded}
\begin{Highlighting}[]
\NormalTok{mcdonalds }\OtherTok{\textless{}{-}}\NormalTok{ fastfood }\SpecialCharTok{\%\textgreater{}\%}
  \FunctionTok{filter}\NormalTok{(restaurant }\SpecialCharTok{==} \StringTok{"Mcdonalds"}\NormalTok{)}
\NormalTok{dairy\_queen }\OtherTok{\textless{}{-}}\NormalTok{ fastfood }\SpecialCharTok{\%\textgreater{}\%}
  \FunctionTok{filter}\NormalTok{(restaurant }\SpecialCharTok{==} \StringTok{"Dairy Queen"}\NormalTok{)}
\end{Highlighting}
\end{Shaded}

\begin{enumerate}
\def\labelenumi{\arabic{enumi}.}
\tightlist
\item
  Make a plot (or plots) to visualize the distributions of the amount of
  calories from fat of the options from these two restaurants. How do
  their centers, shapes, and spreads compare?
\end{enumerate}

Solution:\\
Fat calories in the Mcdonalds food seems to reach higher proportions of
fat with some samples reaching above 1000 fat calories. In the DQ
samples we see a max of around 700 - 750. Both visuals display right
skewed data with a mean around 200.

\begin{Shaded}
\begin{Highlighting}[]
\FunctionTok{ggplot}\NormalTok{(mcdonalds, }\FunctionTok{aes}\NormalTok{(}\AttributeTok{x =}\NormalTok{ cal\_fat)) }\SpecialCharTok{+} 
  \FunctionTok{geom\_histogram}\NormalTok{(}\AttributeTok{binwidth =} \DecValTok{100}\NormalTok{) }\SpecialCharTok{+}
  \FunctionTok{xlab}\NormalTok{(}\StringTok{"Mcdonalds fat calories"}\NormalTok{)}
\end{Highlighting}
\end{Shaded}

\includegraphics[width=0.7\linewidth]{Lab4_normal_distribution_files/figure-latex/unnamed-chunk-2-1}

\begin{Shaded}
\begin{Highlighting}[]
\FunctionTok{ggplot}\NormalTok{(dairy\_queen, }\FunctionTok{aes}\NormalTok{(}\AttributeTok{x =}\NormalTok{ cal\_fat)) }\SpecialCharTok{+}
  \FunctionTok{geom\_histogram}\NormalTok{(}\AttributeTok{binwidth =} \DecValTok{100}\NormalTok{) }\SpecialCharTok{+}
  \FunctionTok{xlab}\NormalTok{(}\StringTok{"Dairy Queen fat calories"}\NormalTok{)}
\end{Highlighting}
\end{Shaded}

\includegraphics[width=0.7\linewidth]{Lab4_normal_distribution_files/figure-latex/unnamed-chunk-2-2}

\subsection{The normal distribution}\label{the-normal-distribution}

In your description of the distributions, did you use words like
\emph{bell-shaped} or \emph{normal}? It's tempting to say so when faced
with a unimodal symmetric distribution.

To see how accurate that description is, you can plot a normal
distribution curve on top of a histogram to see how closely the data
follow a normal distribution. This normal curve should have the same
mean and standard deviation as the data. You'll be focusing on calories
from fat from Dairy Queen products, so let's store them as a separate
object and then calculate some statistics that will be referenced later.

\begin{Shaded}
\begin{Highlighting}[]
\NormalTok{dqmean }\OtherTok{\textless{}{-}} \FunctionTok{mean}\NormalTok{(dairy\_queen}\SpecialCharTok{$}\NormalTok{cal\_fat)}
\NormalTok{dqsd   }\OtherTok{\textless{}{-}} \FunctionTok{sd}\NormalTok{(dairy\_queen}\SpecialCharTok{$}\NormalTok{cal\_fat)}
\end{Highlighting}
\end{Shaded}

Next, you make a density histogram to use as the backdrop and use the
\texttt{lines} function to overlay a normal probability curve. The
difference between a frequency histogram and a density histogram is that
while in a frequency histogram the \emph{heights} of the bars add up to
the total number of observations, in a density histogram the
\emph{areas} of the bars add up to 1. The area of each bar can be
calculated as simply the height \emph{times} the width of the bar. Using
a density histogram allows us to properly overlay a normal distribution
curve over the histogram since the curve is a normal probability density
function that also has area under the curve of 1. Frequency and density
histograms both display the same exact shape; they only differ in their
y-axis. You can verify this by comparing the frequency histogram you
constructed earlier and the density histogram created by the commands
below.

\begin{Shaded}
\begin{Highlighting}[]
\FunctionTok{ggplot}\NormalTok{(}\AttributeTok{data =}\NormalTok{ dairy\_queen, }\FunctionTok{aes}\NormalTok{(}\AttributeTok{x =}\NormalTok{ cal\_fat)) }\SpecialCharTok{+}
        \FunctionTok{geom\_blank}\NormalTok{() }\SpecialCharTok{+}
        \FunctionTok{geom\_histogram}\NormalTok{(}\FunctionTok{aes}\NormalTok{(}\AttributeTok{y =} \FunctionTok{after\_stat}\NormalTok{(density))) }\SpecialCharTok{+}
        \FunctionTok{stat\_function}\NormalTok{(}\AttributeTok{fun =}\NormalTok{ dnorm, }\AttributeTok{args =} \FunctionTok{c}\NormalTok{(}\AttributeTok{mean =}\NormalTok{ dqmean, }\AttributeTok{sd =}\NormalTok{ dqsd), }\AttributeTok{col =} \StringTok{"tomato"}\NormalTok{)}
\end{Highlighting}
\end{Shaded}

\includegraphics[width=0.7\linewidth]{Lab4_normal_distribution_files/figure-latex/hist-height-1}

After initializing a blank plot with \texttt{geom\_blank()}, the
\texttt{ggplot2} package (within the \texttt{tidyverse}) allows us to
add additional layers. The first layer is a density histogram. The
second layer is a statistical function -- the density of the normal
curve, \texttt{dnorm}. We specify that we want the curve to have the
same mean and standard deviation as the column of female heights. The
argument \texttt{col} simply sets the color for the line to be drawn. If
we left it out, the line would be drawn in black.

\begin{enumerate}
\def\labelenumi{\arabic{enumi}.}
\setcounter{enumi}{1}
\tightlist
\item
  Based on the this plot, does it appear that the data follow a nearly
  normal distribution?
\end{enumerate}

\textbf{Response} Yes, with the exception of a few outliers in the
higher end, the data seems to follow nearly normal distribution.

\subsection{Evaluating the normal
distribution}\label{evaluating-the-normal-distribution}

Eyeballing the shape of the histogram is one way to determine if the
data appear to be nearly normally distributed, but it can be frustrating
to decide just how close the histogram is to the curve. An alternative
approach involves constructing a normal probability plot, also called a
normal Q-Q plot for ``quantile-quantile''.

\begin{Shaded}
\begin{Highlighting}[]
\CommentTok{\# Q–Q plot for checking normality}
\FunctionTok{ggplot}\NormalTok{(}\AttributeTok{data =}\NormalTok{ dairy\_queen, }\FunctionTok{aes}\NormalTok{(}\AttributeTok{sample =}\NormalTok{ cal\_fat)) }\SpecialCharTok{+}
  \FunctionTok{geom\_qq}\NormalTok{() }\SpecialCharTok{+}                      \CommentTok{\# plots sample quantiles (points)}
  \FunctionTok{geom\_qq\_line}\NormalTok{() }\SpecialCharTok{+}                 \CommentTok{\# adds the 45° reference line}
  \FunctionTok{labs}\NormalTok{(}
    \AttributeTok{title =} \StringTok{"Q–Q Plot of Calories from Fat"}\NormalTok{,}
    \AttributeTok{x =} \StringTok{"Theoretical Quantiles (Normal)"}\NormalTok{,}
    \AttributeTok{y =} \StringTok{"Sample Quantiles"}
\NormalTok{  )}
\end{Highlighting}
\end{Shaded}

\includegraphics[width=0.7\linewidth]{Lab4_normal_distribution_files/figure-latex/qq-1}

This time, you can use the \texttt{geom\_line()} layer, while specifying
that you will be creating a Q-Q plot with the \texttt{stat} argument.
It's important to note that here, instead of using \texttt{x} in
\texttt{aes()}, you need to use \texttt{sample}.

The x-axis values correspond to the quantiles of a theoretically normal
curve with mean 0 and standard deviation 1 (i.e., the standard normal
distribution). The y-axis values correspond to the quantiles of the
original unstandardized sample data. However, even if we were to
standardize the sample data values, the Q-Q plot would look identical. A
data set that is nearly normal will result in a probability plot where
the points closely follow a diagonal line. Any deviations from normality
leads to deviations of these points from that line.

The plot for Dairy Queen's calories from fat shows points that tend to
follow the line but with some errant points towards the upper tail.
You're left with the same problem that we encountered with the histogram
above: how close is close enough?

A useful way to address this question is to rephrase it as: what do
probability plots look like for data that I \emph{know} came from a
normal distribution? We can answer this by simulating data from a normal
distribution using \texttt{rnorm}.

\begin{Shaded}
\begin{Highlighting}[]
\NormalTok{sim\_norm }\OtherTok{\textless{}{-}} \FunctionTok{rnorm}\NormalTok{(}\AttributeTok{n =} \FunctionTok{nrow}\NormalTok{(dairy\_queen), }\AttributeTok{mean =}\NormalTok{ dqmean, }\AttributeTok{sd =}\NormalTok{ dqsd)}
\end{Highlighting}
\end{Shaded}

The first argument indicates how many numbers you'd like to generate,
which we specify to be the same number of menu items in the
\texttt{dairy\_queen} data set using the \texttt{nrow()} function. The
last two arguments determine the mean and standard deviation of the
normal distribution from which the simulated sample will be generated.
You can take a look at the shape of our simulated data set,
\texttt{sim\_norm}, as well as its normal probability plot.

\begin{enumerate}
\def\labelenumi{\arabic{enumi}.}
\setcounter{enumi}{2}
\tightlist
\item
  Make a normal probability plot of \texttt{sim\_norm}. Do all of the
  points fall on the line? How does this plot compare to the probability
  plot for the real data? (Since \texttt{sim\_norm} is not a dataframe,
  it can be put directly into the \texttt{sample} argument and the
  \texttt{data} argument can be dropped.)
\end{enumerate}

\textbf{Solution} Plotting the sim\_norm data shows data similiar to the
real data but seems to be skewed to the right.

\begin{Shaded}
\begin{Highlighting}[]
\FunctionTok{ggplot}\NormalTok{(}\AttributeTok{data =} \ConstantTok{NULL}\NormalTok{, }\FunctionTok{aes}\NormalTok{(}\AttributeTok{sample =}\NormalTok{ sim\_norm)) }\SpecialCharTok{+}
\FunctionTok{geom\_line}\NormalTok{(}\AttributeTok{stat =} \StringTok{"qq"}\NormalTok{)}
\end{Highlighting}
\end{Shaded}

\includegraphics[width=0.7\linewidth]{Lab4_normal_distribution_files/figure-latex/unnamed-chunk-3-1}

\begin{Shaded}
\begin{Highlighting}[]
\FunctionTok{qqnorm}\NormalTok{(sim\_norm)}
\end{Highlighting}
\end{Shaded}

\includegraphics[width=0.7\linewidth]{Lab4_normal_distribution_files/figure-latex/unnamed-chunk-4-1}

Even better than comparing the original plot to a single plot generated
from a normal distribution is to compare it to many more plots using the
following function. It shows the Q-Q plot corresponding to the original
data in the top left corner, and the Q-Q plots of 8 different simulated
normal data.

\begin{Shaded}
\begin{Highlighting}[]
\CommentTok{\# qqnormsim function definition}
\NormalTok{qqnormsim }\OtherTok{\textless{}{-}} \ControlFlowTok{function}\NormalTok{(sample, data) \{}
\NormalTok{  y }\OtherTok{\textless{}{-}} \FunctionTok{eval}\NormalTok{(}\FunctionTok{substitute}\NormalTok{(sample), data)}
\NormalTok{  simnorm }\OtherTok{\textless{}{-}}\NormalTok{ stats}\SpecialCharTok{::}\FunctionTok{rnorm}\NormalTok{(}
    \AttributeTok{n =} \FunctionTok{length}\NormalTok{(y) }\SpecialCharTok{*} \DecValTok{8}\NormalTok{,}
    \AttributeTok{mean =} \FunctionTok{mean}\NormalTok{(y),}
    \AttributeTok{sd =}\NormalTok{ stats}\SpecialCharTok{::}\FunctionTok{sd}\NormalTok{(y)}
\NormalTok{  )}
  
\NormalTok{  df }\OtherTok{\textless{}{-}} \FunctionTok{data.frame}\NormalTok{(}
    \AttributeTok{x =} \FunctionTok{c}\NormalTok{(y, simnorm),}
    \AttributeTok{plotnum =} \FunctionTok{rep}\NormalTok{(}
      \FunctionTok{c}\NormalTok{(}\StringTok{"data"}\NormalTok{, }\StringTok{"sim 1"}\NormalTok{, }\StringTok{"sim 2"}\NormalTok{, }\StringTok{"sim 3"}\NormalTok{, }\StringTok{"sim 4"}\NormalTok{, }
        \StringTok{"sim 5"}\NormalTok{, }\StringTok{"sim 6"}\NormalTok{, }\StringTok{"sim 7"}\NormalTok{, }\StringTok{"sim 8"}\NormalTok{),}
      \AttributeTok{each =} \FunctionTok{length}\NormalTok{(y)}
\NormalTok{    )}
\NormalTok{  )}
  
\NormalTok{  ggplot2}\SpecialCharTok{::}\FunctionTok{ggplot}\NormalTok{(}\AttributeTok{data =}\NormalTok{ df, ggplot2}\SpecialCharTok{::}\FunctionTok{aes\_string}\NormalTok{(}\AttributeTok{sample =} \StringTok{"x"}\NormalTok{)) }\SpecialCharTok{+}
\NormalTok{    ggplot2}\SpecialCharTok{::}\FunctionTok{stat\_qq}\NormalTok{() }\SpecialCharTok{+}
\NormalTok{    ggplot2}\SpecialCharTok{::}\FunctionTok{geom\_qq\_line}\NormalTok{(}\AttributeTok{color =} \StringTok{"darkred"}\NormalTok{, }\AttributeTok{linetype =} \StringTok{"dashed"}\NormalTok{) }\SpecialCharTok{+}  \CommentTok{\# added reference line}
\NormalTok{    ggplot2}\SpecialCharTok{::}\FunctionTok{facet\_wrap}\NormalTok{(}\SpecialCharTok{\textasciitilde{}}\NormalTok{plotnum) }\SpecialCharTok{+}
\NormalTok{    ggplot2}\SpecialCharTok{::}\FunctionTok{labs}\NormalTok{(}
      \AttributeTok{title =} \StringTok{"Q–Q Plots: Data vs. Simulated Normal Samples"}\NormalTok{,}
      \AttributeTok{x =} \StringTok{"Theoretical Quantiles (Normal)"}\NormalTok{,}
      \AttributeTok{y =} \StringTok{"Sample Quantiles"}
\NormalTok{    ) }
\NormalTok{\}}
\FunctionTok{set.seed}\NormalTok{(}\DecValTok{21}\NormalTok{)}
\FunctionTok{qqnormsim}\NormalTok{(}\AttributeTok{sample =}\NormalTok{ cal\_fat, }\AttributeTok{data =}\NormalTok{ dairy\_queen)}
\end{Highlighting}
\end{Shaded}

\includegraphics[width=0.7\linewidth]{Lab4_normal_distribution_files/figure-latex/qqnormsim-1}

\begin{enumerate}
\def\labelenumi{\arabic{enumi}.}
\setcounter{enumi}{3}
\tightlist
\item
  Does the normal probability plot for the calories from fat look
  similar to the plots created for the simulated data? That is, do the
  plots provide evidence that the calories from fat are nearly normal?
\end{enumerate}

\textbf{Solution}\\
Yes the simulated normal probability plots looks very much like the data
plot for fat calories.

\begin{enumerate}
\def\labelenumi{\arabic{enumi}.}
\setcounter{enumi}{4}
\tightlist
\item
  Using the same technique, determine whether or not the calories from
  McDonald's menu appear to come from a normal distribution.
\end{enumerate}

\textbf{Solution}\\
We can see that the data for Mcdonalds is close to the simulated normal
distributions however it begins to diverge towards the right end going
higher due to the samples containing larger protoin of fat calories.
This is showing that the Mcdonalds data is right skewed.

\begin{Shaded}
\begin{Highlighting}[]
\FunctionTok{qqnormsim}\NormalTok{(}\AttributeTok{sample =}\NormalTok{ cal\_fat, }\AttributeTok{data =}\NormalTok{ mcdonalds)}
\end{Highlighting}
\end{Shaded}

\includegraphics[width=0.7\linewidth]{Lab4_normal_distribution_files/figure-latex/unnamed-chunk-5-1}

\subsection{Normal probabilities}\label{normal-probabilities}

Okay, so now you have a slew of tools to judge whether or not a variable
is normally distributed. Why should you care?

It turns out that statisticians know a lot about the normal
distribution. Once you decide that a random variable is approximately
normal, you can answer all sorts of questions about that variable
related to probability. Take, for example, the question of, ``What is
the probability that a randomly chosen Dairy Queen product has more than
600 calories from fat?''

If we assume that the calories from fat from Dairy Queen's menu are
normally distributed (a very close approximation is also okay), we can
find this probability by calculating a Z score and consulting a Z table
(also called a normal probability table). In R, this is done in one step
with the function \texttt{pnorm()}.

\begin{Shaded}
\begin{Highlighting}[]
\DecValTok{1} \SpecialCharTok{{-}} \FunctionTok{pnorm}\NormalTok{(}\AttributeTok{q =} \DecValTok{600}\NormalTok{, }\AttributeTok{mean =}\NormalTok{ dqmean, }\AttributeTok{sd =}\NormalTok{ dqsd)}
\end{Highlighting}
\end{Shaded}

\begin{verbatim}
## [1] 0.01501523
\end{verbatim}

Note that the function \texttt{pnorm()} gives the area under the normal
curve below a given value, \texttt{q}, with a given mean and standard
deviation. Since we're interested in the probability that a Dairy Queen
item has more than 600 calories from fat, we have to take one minus that
probability.

Assuming a normal distribution has allowed us to calculate a theoretical
probability. If we want to calculate the probability empirically, we
simply need to determine how many observations fall above 600 then
divide this number by the total sample size.

\begin{Shaded}
\begin{Highlighting}[]
\NormalTok{dairy\_queen }\SpecialCharTok{\%\textgreater{}\%} 
  \FunctionTok{filter}\NormalTok{(cal\_fat }\SpecialCharTok{\textgreater{}} \DecValTok{600}\NormalTok{) }\SpecialCharTok{\%\textgreater{}\%}
  \FunctionTok{summarise}\NormalTok{(}\AttributeTok{percent =} \FunctionTok{n}\NormalTok{() }\SpecialCharTok{/} \FunctionTok{nrow}\NormalTok{(dairy\_queen))}
\end{Highlighting}
\end{Shaded}

\begin{verbatim}
## # A tibble: 1 x 1
##   percent
##     <dbl>
## 1  0.0476
\end{verbatim}

Although the probabilities are not exactly the same, they are reasonably
close. The closer that your distribution is to being normal, the more
accurate the theoretical probabilities will be.

\begin{enumerate}
\def\labelenumi{\arabic{enumi}.}
\setcounter{enumi}{5}
\tightlist
\item
  Write out two probability questions that you would like to answer
  about any of the restaurants in this dataset. Calculate those
  probabilities using both the theoretical normal distribution as well
  as the empirical distribution (four probabilities in all). Which one
  had a closer agreement between the two methods?
\end{enumerate}

\textbf{Solution} First lets calculate the probability that a randomly
chosen item from Mcdonalds has more than 500 calories.

\begin{Shaded}
\begin{Highlighting}[]
\NormalTok{mcdonalds\_mean }\OtherTok{\textless{}{-}} \FunctionTok{mean}\NormalTok{(mcdonalds}\SpecialCharTok{$}\NormalTok{cal\_fat)}
\NormalTok{mcdonalds\_sd   }\OtherTok{\textless{}{-}} \FunctionTok{sd}\NormalTok{(mcdonalds}\SpecialCharTok{$}\NormalTok{cal\_fat)}
\FunctionTok{pnorm}\NormalTok{(}\DecValTok{500}\NormalTok{, mcdonalds\_mean, mcdonalds\_sd)}
\end{Highlighting}
\end{Shaded}

\begin{verbatim}
## [1] 0.834105
\end{verbatim}

\begin{center}\rule{0.5\linewidth}{0.5pt}\end{center}

\subsection{More Practice}\label{more-practice}

\begin{enumerate}
\def\labelenumi{\arabic{enumi}.}
\setcounter{enumi}{6}
\tightlist
\item
  Now let's consider some of the other variables in the dataset. Out of
  all the different restaurants, which ones' distribution is the closest
  to normal for sodium?
\end{enumerate}

\textbf{Solution} Burger King seems to have the closest distribution to
normal for sodium.

\begin{Shaded}
\begin{Highlighting}[]
\NormalTok{normal\_sodium }\OtherTok{\textless{}{-}}\NormalTok{ fastfood }\SpecialCharTok{\%\textgreater{}\%}
  \FunctionTok{group\_by}\NormalTok{(restaurant) }\SpecialCharTok{\%\textgreater{}\%}
  \FunctionTok{summarise}\NormalTok{(}\AttributeTok{shapiro\_p\_value =} \FunctionTok{shapiro.test}\NormalTok{(sodium)}\SpecialCharTok{$}\NormalTok{p.value)}

\NormalTok{normal\_restaurants }\OtherTok{\textless{}{-}}\NormalTok{ normal\_sodium }\SpecialCharTok{\%\textgreater{}\%}
  \FunctionTok{filter}\NormalTok{(shapiro\_p\_value }\SpecialCharTok{\textgreater{}} \FloatTok{0.05}\NormalTok{)}

\FunctionTok{print}\NormalTok{(normal\_restaurants)}
\end{Highlighting}
\end{Shaded}

\begin{verbatim}
## # A tibble: 2 x 2
##   restaurant  shapiro_p_value
##   <chr>                 <dbl>
## 1 Arbys                 0.199
## 2 Burger King           0.133
\end{verbatim}

\begin{enumerate}
\def\labelenumi{\arabic{enumi}.}
\setcounter{enumi}{7}
\tightlist
\item
  Note that some of the normal probability plots for sodium
  distributions seem to have a stepwise pattern. why do you think this
  might be the case?
\end{enumerate}

\textbf{Solution} The data is likely discrete and not continuous, this
causes a step like pattern on a normal probability plot.

\begin{center}\rule{0.5\linewidth}{0.5pt}\end{center}

\pandocbounded{\includegraphics[keepaspectratio]{https://i.creativecommons.org/l/by-sa/4.0/88x31.png}}This
work is licensed under a Creative Commons Attribution-ShareAlike 4.0
International License.

\end{document}
